\section{Числовий аналіз і обговорення результатів}
\label{sec:numerical-results}

Числовий аналіз виконано для нормованих параметрів
$\,d_i/l$, $\delta_i'$ і $J_i'$, визначених у
розділі~\ref{sec:derivation}. Фізично допустимими вважаємо лише ті
комбінації матеріалу, кута та знака навантаження, для яких одночасно
виконуються умови
\begin{equation*}
  Cg_2<0,\qquad CQ_i>0.
\end{equation*}
Перша нерівність забезпечує стискання берегів міжфазної зсувної тріщини,
а друга --- відривне нормальне напруження на бісектрисі кутового сектора матеріалу, в якому
утворюється зона передруйнування.

Для базової пари матеріалів
\begin{equation*}
  E_1/E_2=0.5,\qquad \nu_1=\nu_2=0.3,
\end{equation*}
зміна знака множника $g_2$ відбувається при
\begin{equation*}
  \alpha_1=12.9202^\circ,\qquad
  \alpha_2=107.0897^\circ.
\end{equation*}
Ці значення отримано безпосереднім розв'язуванням умов
$D_0(-1-\lambda_0)=0$ і $g_2=0$. За $C>0$ зона передруйнування є фізично
допустимою в матеріалі~2 при $0<\alpha<\alpha_1$ і в матеріалі~1 при
$\alpha_2<\alpha<180^\circ$. За $C<0$ відповідні області мають вигляд
$\alpha_1<\alpha<90^\circ$ для матеріалу~1 та
$90^\circ<\alpha<\alpha_2$ для матеріалу~2. Значення
$\alpha=90^\circ$ відповідає прямолінійній межі поділу і надалі
розглядається як вироджена аналітична границя, а не як звичайна внутрішня
точка числової гілки.

Для оцінки впливу пружного контрасту ті самі два переходи обчислено для
низки значень $E_1/E_2<1$ при $\nu_1=\nu_2=0.3$. Результати наведено в
табл.~\ref{tab:g2-transition-angles}. Зі зростанням $E_1/E_2$ перший
перехідний кут $\alpha_1$ монотонно зростає, тоді як $\alpha_2$
монотонно зменшується, тобто пружний контраст систематично змінює межі
кутових інтервалів фізичної допустимості.

\begin{table}[htbp]
  \centering
  \caption{Кути зміни знака множника $g_2$ при $\nu_1=\nu_2=0.3$.}
  \label{tab:g2-transition-angles}
  \small
  \setlength{\tabcolsep}{3.4pt}
  \begin{tabular}{@{}lrrrrrrrrrrr@{}}
    \toprule
    $E_1/E_2$ & 0.01 & 0.10 & 0.20 & 0.30 & 0.40 & 0.50 & 0.60 & 0.70 & 0.80 & 0.90 & 0.99 \\
    \midrule
    $\alpha_1$, $^\circ$ & 2.58 & 7.03 & 9.25 & 10.77 & 11.95 & 12.92 & 13.74 & 14.45 & 15.08 & 15.65 & 16.11 \\
    $\alpha_2$, $^\circ$ & 108.02 & 107.90 & 107.72 & 107.52 & 107.30 & 107.09 & 106.88 & 106.69 & 106.50 & 106.32 & 106.17 \\
    \bottomrule
  \end{tabular}
\end{table}


Випадки $E_1/E_2>1$ не є незалежними: вони відновлюються перестановкою
матеріалів
\begin{equation*}
  (E_1,\nu_1)\leftrightarrow(E_2,\nu_2),
  \qquad
  \alpha\leftrightarrow\pi-\alpha,
  \qquad
  1\leftrightarrow2.
\end{equation*}
За $\nu_1=\nu_2$ це, зокрема, відповідає заміні
$E_1/E_2\leftrightarrow E_2/E_1$, тому в табл.~\ref{tab:g2-transition-angles}
достатньо навести лише $E_1/E_2<1$.

\subsection{Залежність від нормованого навантаження}
\label{subsec:figure2}

На рис.~\ref{fig:load-dependence} наведено залежності параметрів зони
передруйнування від
\begin{equation*}
  \sigma'=\frac{Cl^{\lambda_0}}{\sigma_i}
\end{equation*}
для чотирьох контрольних конфігурацій:
$\alpha=10^\circ$ ($i=2$, $C>0$),
$45^\circ$ ($i=1$, $C<0$),
$105^\circ$ ($i=2$, $C<0$) і
$135^\circ$ ($i=1$, $C>0$). Для кожної конфігурації показано лише
фізично допустиму піввісь навантаження.

\resultfigure{figure2_recalculated.pdf}{%
Залежність нормованих параметрів зони передруйнування від
безрозмірного навантаження $\sigma'$ для чотирьох контрольних кутів і
фізично допустимих матеріалів.}{fig:load-dependence}

У всіх чотирьох випадках $d_i/l$, $\delta_i'$ і $J_i'$ монотонно
зростають зі збільшенням $|\sigma'|$. Така синхронність безпосередньо
випливає з аналітичних співвідношень
$\delta_i'\propto d_i/l$ і $J_i'\propto d_i/l$ при фіксованій
геометрії. Серед наведених контрольних конфігурацій найбільші значення
усіх трьох параметрів при $|\sigma'|=0.5$ отримано для
$\alpha=45^\circ$, $C<0$; зокрема,
\begin{equation*}
  \frac{d_1}{l}=0.14743436,\qquad
  \delta_1'=0.14578703,\qquad
  J_1'=0.08362039.
\end{equation*}
Цей результат важливий також для перевірки маломасштабного припущення:
збільшення навантаження одночасно збільшує локальні деформаційний та
енергетичний параметри і наближає модель до межі її асимптотичної
застосовності.

Для фіксованої конфігурації розкриття $\delta_i$ та зональний
енергетичний параметр $J_i$ однозначно пов'язані з рівноважною довжиною
$d_i$. Тому графіки $\delta_i'$ і $J_i'$ можуть бути використані як
деформаційне та енергетичне представлення тієї самої однопараметричної
сім'ї локальних станів. Критерії $\delta_i=\delta_{i\mathrm{c}}$ та
$J_i=J_{i\mathrm{c}}$ є еквівалентними лише за узгодженого вибору
незалежно заданих критичних матеріальних параметрів
$\delta_{i\mathrm{c}}$ і $J_{i\mathrm{c}}$.

\subsection{Залежність від кута зламу межі поділу}
\label{subsec:figure3}

Рис.~\ref{fig:angle-dependence} показує залежність нормованих параметрів
зони від півкута $\alpha$ для тієї самої базової пари матеріалів при
$\sigma'=\pm0.5$. На рисунку залишено лише комбінації матеріалу та знака
навантаження, що задовольняють умовам фізичної допустимості. Переходи між
допустимими гілками відбуваються при $\alpha_1$,
$90^\circ$ і $\alpha_2$.

\resultfigure{figure3_recalculated.pdf}{%
Залежність нормованих параметрів зони передруйнування від півкута
$\alpha$ при $\sigma'=\pm0.5$. Показано лише фізично допустимі комбінації
матеріалу та знака навантаження; $\alpha=90^\circ$ є аналітичною межею.}%
{fig:angle-dependence}

На двох широких допустимих гілках матеріалу~1 кутові залежності мають
виражені внутрішні максимуми, тоді як на коротких гілках матеріалу~2
найбільші значення досягаються поблизу меж допустимості. Після уточнення
аналітичних формул максимуми $d_i/l$, $\delta_i'$ і $J_i'$ є близькими,
але не збігаються за кутом. Для синьої гілки матеріалу~1 при $C<0$
одноградусний розрахунок дає
\begin{equation*}
\begin{aligned}
  \max(d_1/l)&=0.17793332 &&\text{при }\alpha=35^\circ,\\
  \max\delta_1'&=0.20770529 &&\text{при }\alpha=32^\circ,\\
  \max J_1'&=0.11391440 &&\text{при }\alpha=33^\circ.
\end{aligned}
\end{equation*}
Ця близькість узгоджується з тим, що при кожному фіксованому $\alpha$
величини $\delta_i'$ і $J_i'$ пропорційні $d_i/l$, але самі коефіцієнти
пропорційності залежать від кута. Отже, рис.~\ref{fig:angle-dependence}
характеризує вплив геометрії на досягнуті за заданого навантаження
деформаційний та енергетичний параметри. Для переходу від цих кривих до
кутової залежності критичного навантаження необхідно задати відповідний
критичний матеріальний параметр і для кожного $\alpha$ окремо розв'язати
рівняння критерію.

Максимум синьої гілки на рис.~\ref{fig:angle-dependence} відповідає
$d_i/l$ порядку $0.18$. Це значення вже не є беззастережно малим порівняно
з одиницею, тоді як побудова локальної задачі ґрунтується на припущенні
$d_i/l\ll1$. Тому повна крива корисна для аналізу структури розв'язку,
але кількісну інтерпретацію області поблизу найбільшого максимуму слід
проводити з урахуванням цієї межі застосовності.

Числове обчислення плюс-факторів для рис.~\ref{fig:angle-dependence}
потребує особливої уваги при $\alpha\to0$ та
$\alpha\to180^\circ$. Фіксований контур факторизації повільно збігається
біля цих меж. У розрахунках застосовано адаптивне розширення контуру; при
додатковому збільшенні хвоста максимальна зміна перевірених
плюс-факторів становила $8.69\times10^{-13}$. Тому
рис.~\ref{fig:angle-dependence} побудовано за збіжними адаптивними
значеннями.

\subsection{Показники сингулярності та фізична допустимість}
\label{subsec:figure4}

Для порівняння розглянемо три локальні конфігурації: інтактну
біматеріальну кутову точку, кутову точку із заданою стаціонарною
міжфазною зсувною тріщиною та конфігурацію тріщина--зона
передруйнування. Для базової пари матеріалів їхні сингулярності
характеризуються показниками $\lambda_0$, $\lambda$ та $\lambda_i$
відповідно.

\resultfigure{figure4_admissible_segments.pdf}{%
Показники сингулярності та фізично допустимі сегменти для базової пари
матеріалів. Повні математичні гілки та фізично допустимі ділянки
розрізняються відповідно до умов $Cg_2<0$ і $CQ_i>0$.}%
{fig:singularity-exponents}

У межах фізично допустимих сегментів для дослідженої базової пари
матеріалів конфігурація з міжфазною зсувною тріщиною має істотно
сильнішу сингулярність, ніж інтактна кутова точка: модуль $\lambda$
перевищує модуль $\lambda_0$. Додавання зони передруйнування частково
послаблює цю концентрацію, причому для відповідного матеріалу виконується
співвідношення
\begin{equation*}
  |\lambda_0|<|\lambda_i|<|\lambda|.
\end{equation*}
Отже, зона передруйнування виконує локальну екранувальну роль: вона
знижує порядок сингулярності, спричинений міжфазною зсувною тріщиною,
але не усуває концентрацію повністю. Це співвідношення тут
встановлюється для розглянутої пари матеріалів і фізично допустимих
гілок; його не слід без додаткової перевірки переносити на довільні
пружні контрасти.

Сильніша сингулярність конфігурації з міжфазною тріщиною створює
передумови для локального руйнування одного з матеріалів, однак саме існування зони
визначається не лише величиною показника сингулярності, а й умовами
$Cg_2<0$, $CQ_i>0$, опором відриву $\sigma_i$ та рівнем зовнішнього
навантаження. Така інтерпретація пов'язує послідовність
$\lambda_0\to\lambda\to\lambda_i$ із механічною моделлю початкової стадії
відгалуження, не надаючи самій сингулярності статусу достатнього критерію
руйнування.
